% ============================================================
% config.tex — EJEMPLO LLENO (proyecto Braille)
% ============================================================

\newcommand{\Titulo}{Diseño y evaluación de una red neuronal convolucional ligera para el reconocimiento óptico de Braille Grado 1 en dispositivos móviles: efecto en la retroalimentación docente en aulas inclusivas de Moquegua e Ilo}
\newcommand{\AutorA}{Rony Frank Churata Huania}
\newcommand{\AutorB}{Erick Jean Ramos Vera}
\newcommand{\Filiacion}{Escuela Profesional de Ingeniería de Sistemas e Informática, Universidad Nacional de Moquegua}
\newcommand{\Curso}{Trabajo de Investigación}
\newcommand{\Docente}{Ing. Raul Donato Cruz Huacac}
\newcommand{\Fecha}{8 de octubre de 2026}

\newcommand{\PalabrasClave}{reconocimiento óptico de Braille; red neuronal convolucional ligera; detección de objetos; evaluación de modelos; educación inclusiva}
\newcommand{\Resumen}{%
En las aulas de Educación Básica Regular de Moquegua e Ilo, los estudiantes con
ceguera escriben en Braille Grado 1 con regleta y punzón sobre cartulina, un
relieve sin contraste cromático que el docente de aula no puede leer, por lo que
la retroalimentación formativa queda diferida hasta la visita del especialista.
Este proyecto diseña y evalúa una red neuronal convolucional ligera para el
reconocimiento óptico de Braille Grado 1 en español, ejecutable sin conexión en
teléfonos de gama media/baja. Se plantean tres objetivos: comparar la precisión
y la robustez de configuraciones basadas en YOLOv8n y YOLOv11 (FP32 e INT8);
determinar el efecto de la cuantización INT8 sobre la eficiencia computacional; y
evaluar el efecto del sistema en el tiempo de revisión docente. El desempeño se
mide con mAP@0.5, F1 macro, Kappa de Cohen, AUC-ROC macro, CER, WER, latencia,
memoria y tamaño del modelo, sobre validación cruzada agrupada por página con un
corpus de 40 páginas escolares (480 fotografías). El tercer objetivo se aborda
con un cuasiexperimento de medidas repetidas con 30 docentes de las UGEL Ilo y
Mariscal Nieto. Se espera identificar la configuración con mejor compromiso
precisión--eficiencia--robustez y aportar un corpus anotado de Braille escolar
manuscrito en español peruano.
}

\PortadaProfesionalfalse
\newcommand{\RunningHead}{RECONOCIMIENTO ÓPTICO DE BRAILLE CON CNN LIGERA}
\ConIndicetrue
